\chapter{Conclusions and Forward Work} \label{ch:conclusions}

This dissertation establishes dust-enabled just-in-time collision avoidance as a coupled architecture-under-uncertainty problem. Hazardous-event evidence, transfer physics, dust-cloud uncertainty, and optimizer comparison are evaluated within one traceable workflow rather than as separate analyses.

\section{Core Contributions} \label{se:conclusions_contributions}

The dissertation contributes four linked methodological pieces. First, it builds a catalogue-first unsafe-event evidence pipeline so candidate constellations are judged against reusable hazardous conjunction opportunities rather than disconnected synthetic examples. Second, it formulates dust-enabled \gls{jca} as a deployer-architecture problem, connecting fixed-epoch transfer replay, release control, deterministic intercepted-mass checks, raw-cloud capture sizing, and mission scoring in one traceable chain. Third, it develops a layered uncertainty-propagation and scoring framework, enabled by a custom multi-method Rust propagator using Battin's Encke formulation, that separates adaptive Beta/\gls{hpd} search diagnostics from exact canonical authority: the canonical \(N=64\) replay bank prices publication coverage at the exact \(q=0.9\) quantile, that is 58 of 64 unsafe events, and retains multiple maneuver--mass solutions on that lane, while search-lane rows are reported as single maneuver--mass pairs by construction because the sealed coverage-front band is \texttt{Qualified}, which no search row reaches under the sealed posterior stop rule. Fourth, it defines a shared-evidence descriptive framework across three architecture families and six optimizer methods while prohibiting global and per-family optimizer winners. Current attainment and operational risk-reduction claims remain held until evidence sealed under the current compiled science authority and receipt chain is independently reconstructed.

These contributions are independent of any one presentation of pairwise optimizer diagnostics or topology counts. The empirical conclusions attach to the reported evidence package, while the methodological contribution is the framework that makes those conclusions readable.

\section{Empirical Closure Points} \label{se:conclusions_empirical_to_finalize}

The empirical close of the dissertation rests on three questions:

\begin{enumerate}[leftmargin=*]
    \item what descriptive differences among optimizer methods are supported by realized-work and final-front evidence from a single-realization campaign, why the reported campaign shapes cannot support an interval or directional-agreement statement, and what independent-replicate shape would be needed to obtain one;
    \item which deployer architecture patterns are supported by the Part~B constellation read, and how strongly those patterns depend on catalogue scope, family encoding, and solver/scoring settings;
    \item which sensitivity findings are robust enough to state as dissertation conclusions rather than as diagnostic or follow-on evidence.
\end{enumerate}

This framing keeps pairwise optimizer diagnostics, exact front counts, architecture sufficiency percentages, and sensitivity margins tied to the evidence that produces them.

\section{Limitations and Boundary Conditions} \label{se:conclusions_limits}

The manuscript carries explicit limits. Comparative statements remain tied to the hazardous-event bank, mission-scoring semantics, family encodings, solver settings, and campaign shape used for the reported evidence package. Fidelity and uncertainty-transport choices remain important near acceptance boundaries, so claims stay attached to the reported baseline rather than drifting into model-form generalizations. Constellation-level conclusions remain scoped to the families studied here; they do not exhaust every possible remediation architecture. All three evaluated families are analytically structured, so whether an unstructured, plane-constrained architecture could outperform them remains open: that encoding was scoped for this study but is not implemented in the reported evaluation stack, and no evidence here bounds it in either direction.

Those limits tell the reader how far the conclusions travel. The dissertation defends its framework and problem formulation without turning campaign-specific margins into universal statements.

\section{Forward Work} \label{se:conclusions_remaining_work}

The natural next layer is broader empirical stress testing. Extensions include richer policy coupling, additional catalogue epochs, broader traffic-management implications, deeper force-model sensitivity sweeps, and further validation of the dust-propagation representation. These extensions build on the same evidence chain rather than replacing it.

Several pieces of that chain are implemented and sealed but not reported empirically here, and they bound what a follow-on campaign must re-measure rather than re-build: the closed-form minimum-separation certificate that replaced the sampled constellation screen, the regenerated representative event banks and their byte-identical generation receipt, the shared-altitude fairness envelope now enforced on the Flower encoding, and the post-Hybrid verification and validation lane. A campaign flown on that stack supersedes every empirical value quoted in \cref{ch:results}.

\section{Closing Perspective} \label{se:conclusions_closing}

Taken as a methodological whole, the dissertation gives dust-enabled just-in-time collision avoidance a concrete and conservative design form. It evaluates deployer architectures through hazardous-event evidence, transfer feasibility, dust-cloud uncertainty, and shared optimizer controls rather than through isolated coverage or average-case surrogates. That discipline is the through-line from problem formulation to empirical interpretation.
